\scanpage{5}\mspage{2}
\begingroup
\makeatletter\tagsleft@true\makeatother
2. \uline{Change} \uline{of} \uline{coordinates}. Suppose $E':\quad y^{\prime\,2}+a_1'x'y'+\ldotp\,\ldotp\,\ldotp\,\ldotp$ is another curve of the same type as $E$, and $f:E'\overset{\sim}{\longrightarrow}E$ an isomorphism over $K$. carrying $O'$ into $O$. Then there are $r,s,t$ and $u\ne0$ in $K$ such that
\[
\begin{aligned}
x\circ f&=u^2x'+r & \omega\circ f&=u^{-1}\omega'.\\
y\circ f&=u^3y'+su^2x'+t
\end{aligned}
\tag*{(2.1)}
\]

The coefficients $a_i'$ are related to the $a_i$ \deleted{\illegible{1}} and the $b_i'$ to the $b_i$ by the formulas:
\begingroup\scriptsize
\[
\begin{aligned}
ua_1'&=a_1+2s & u^8b_8'&=b_8+3rb_6+3r^2b_4+r^3b_2+3r^{\unsure{4}}\\
u^2a_2'&=a_2-sa_1+3r-s^2 & u^6b_6'&=b_6+2rb_4+r^2b_2+4r^3\\
u^3a_3'&=a_3+ra_1+2t=F_y(r,t) & u^4b_4'&=b_4+rb_2+6r^2\\
u^4a_4'&=a_4-sa_3+2ra_2-(t+rs)a_1+3r^2-2st
 & u^2b_2'&=b_2+12r.\\
 &= -F_x(r,t)-sF_y(r,t)\\
u^6a_6'&=a_6+ra_4+r^2a_2+r^3-ta_3-t^2-rta_1\\
 &= -F(r,t)
\end{aligned}
\tag*{(2.2)}
\]
\endgroup

From these we check that:
\[
u^4c_4'=c_4\qquad u^6c_6'=c_6\qquad
u^{12}\Delta'=\Delta\qquad j'=j.
\tag*{(2.3)}
\]

\uline{Example} \unsure{`|``}generic $E_j$''. The equation
\[
y^2+xy=x^3-\frac{36}{j-1728}x-\frac{1}{j-1728}
\tag*{(2.4)}
\]
has $c_4=c_6=\dfrac{j}{j-1728}$ and $\Delta=\dfrac{j^2}{(j-1728)^3}$. Hence for $j\ne0,1728$ it gives a curve with invariant '$j$.

Using the formulas above, it is easy to show that $j$ can be arbitrary in $K$, and that for $K$ algebraically closed, two $E$'s with the same $j$ are \unsure{isomorphic} It is also easy to compute the group of automorphisms (but not, alas, the ring of endomorphisms). See Serre's ``Courbes Elliptiques - Formulaire'' for the details.
\endgroup

% archive sheet number: 4
